\begin{table}[!h]
\centering
\caption{\label{tab:resumo_socioespacial}Medidas-resumo das variáveis socioespaciais por Grande Região (Censo 2022).}
\centering
\resizebox{\ifdim\width>\linewidth\linewidth\else\width\fi}{!}{
\fontsize{9}{11}\selectfont
\begin{tabular}[t]{lcccccc}
\toprule
Medida & Brasil & Norte & Nordeste & Sudeste & Sul & Centro-Oeste\\
\midrule
\addlinespace[0.3em]
\multicolumn{7}{l}{\textbf{Abastecimento de água por rede geral (\%)}}\\
\hspace{1em}Média (D.P.) & 71,7 (19,6) & 53,4 (22,8) & 67,3 (20,0) & 79,1 (15,6) & 74,3 (18,1) & 73,5 (15,4)\\
\hspace{1em}Mediana (IIQ) & 75,8 (26,2) & 57,6 (32,9) & 70,5 (25,4) & 82,8 (21,7) & 78,7 (25,8) & 75,7 (17,8)\\
\hspace{1em}Mínimo–Máximo & 0,0 – 100,0 & 0,0 – 95,0 & 0,0 – 99,6 & 22,5 – 100,0 & 10,2 – 99,6 & 3,4 – 96,6\\
\addlinespace[0.3em]
\multicolumn{7}{l}{\textbf{Esgotamento sanitário por rede (\%)}}\\
\hspace{1em}Média (D.P.) & 37,7 (33,0) & 7,4 (12,4) & 24,4 (23,9) & 70,6 (23,4) & 29,8 (29,3) & 20,7 (25,4)\\
\hspace{1em}Mediana (IIQ) & 32,3 (63,7) & 2,5 (6,4) & 17,0 (38,1) & 76,8 (32,3) & 18,9 (51,3) & 6,6 (34,0)\\
\hspace{1em}Mínimo–Máximo & 0,0 – 99,9 & 0,0 – 76,6 & 0,0 – 95,1 & 0,1 – 99,9 & 0,0 – 98,1 & 0,0 – 94,3\\
\addlinespace[0.3em]
\multicolumn{7}{l}{\textbf{Coleta de lixo (\%)}}\\
\hspace{1em}Média (D.P.) & 80,2 (17,4) & 66,2 (17,0) & 70,7 (17,8) & 89,1 (12,5) & 86,8 (13,5) & 81,6 (13,0)\\
\hspace{1em}Mediana (IIQ) & 85,1 (25,5) & 66,7 (25,8) & 72,5 (25,1) & 94,2 (12,3) & 91,3 (16,4) & 84,6 (17,1)\\
\hspace{1em}Mínimo–Máximo & 1,7 – 100,0 & 13,6 – 98,8 & 1,7 – 99,6 & 24,1 – 100,0 & 31,3 – 100,0 & 29,1 – 99,8\\
\addlinespace[0.3em]
\multicolumn{7}{l}{\textbf{População em favelas (\%)}}\\
\hspace{1em}Média (D.P.) & 1,0 (4,3) & 3,1 (8,9) & 0,8 (3,9) & 1,2 (4,2) & 0,5 (2,5) & 0,2 (1,1)\\
\hspace{1em}Mediana (IIQ) & 0,0 (0,0) & 0,0 (0,0) & 0,0 (0,0) & 0,0 (0,0) & 0,0 (0,0) & 0,0 (0,0)\\
\hspace{1em}Mínimo–Máximo & 0,0 – 69,2 & 0,0 – 69,2 & 0,0 – 45,4 & 0,0 – 42,6 & 0,0 – 47,2 & 0,0 – 11,8\\
\addlinespace[0.3em]
\multicolumn{7}{l}{\textbf{Densidade demográfica (hab/km²)}}\\
\hspace{1em}Média (D.P.) & 113,8 (585,8) & 23,1 (142,2) & 89,1 (395,6) & 205,7 (942,0) & 87,0 (285,8) & 36,1 (207,1)\\
\hspace{1em}Mediana (IIQ) & 24,1 (41,8) & 4,1 (9,0) & 28,6 (48,5) & 30,8 (53,0) & 25,5 (33,9) & 5,8 (12,2)\\
\hspace{1em}Mínimo–Máximo & 0,2 – 13.328,0 & 0,2 – 2.501,1 & 1,0 – 7.958,5 & 1,2 – 13.328,0 & 1,5 – 4.065,9 & 0,3 – 3.221,8\\
\bottomrule
\end{tabular}}
\end{table}
